This element models Compton back-scattering of a counter-propagating (or crossing-angle) laser pulse off the electron
beam, following and extending Pan {\em et al.}  \cite{PhysRevAccelBeams.22.040702}.  It is a zero-length, per-particle
stochastic element: on each pass every particle has a probability of undergoing a single Compton scatter, and if it does
its transverse angles and fractional momentum deviation are changed according to Eq.~(4) of that reference.  The element
may optionally record the emitted (back-scattered) photons to an SDDS file for $\gamma$-source spectrum studies.

\paragraph{Recoil model.}
Let $\gamma$ and $\beta$ be the Lorentz factor and velocity of the {\em individual}
electron (computed from its own momentum $P=(1+\delta)P_0$, not the reference
value, so a beam with energy spread scatters with the correct per-electron
kinematics), $E_e=\gamma m_e c^2$ its energy, $E_h = h c/\lambda$ the lab-frame
laser-photon energy (set by \verb|LASER_WAVELENGTH|), and $\theta_0$ the collision
angle between the electron and laser directions (\verb|COLLISION_ANGLE|;
$\theta_0=\pi$ is head-on).  Define the rest-frame incident photon energy
\begin{equation}
E' = \gamma E_h (1 - \beta\cos\theta_0).
\end{equation}
For a scatter with photon polar angle $\theta$ (measured from the electron
direction) and azimuth $\phi$, the coordinate changes applied are
\begin{eqnarray}
x' &=& x'_0 - \frac{E'}{E_e\beta}\sin\theta\cos\phi, \\
y' &=& y'_0 - \frac{E'}{E_e\beta}\sin\theta\sin\phi, \\
\delta &=& \delta_0 + \frac{\gamma^2 E_h}{E_e\beta}(\cos\theta_0-\beta)
                    - \frac{\gamma E'}{E_e\beta}\cos\theta .
\end{eqnarray}
The change in $\delta$ is applied so as to preserve the particle arrival time.
The maximum energy transferred to the photon (the Compton edge) is recovered at
$\theta=0$ and is $E_{\gamma,\mathrm{max}}\approx 4\gamma^2 E_h$ for a head-on
collision.

The Eq.~(4) mapping above is a linearized lab-frame recoil: the transverse terms
are the paraxial projection $x'=p_x/p_z$, and the energy change is evaluated using
the {\em incident} rest-frame photon energy $E'$.  It is therefore exact only in
the Thomson limit $\epsilon\equiv E'/(m_e c^2)\to 0$; for $\epsilon$ not small it
overestimates the electron energy loss because it omits the rest-frame Compton
down-shift.  Setting \verb|EXACT_RECOIL|$=1$ replaces this mapping with an exact,
per-particle treatment: the electron and incident photon four-momenta are boosted
into the electron rest frame (boost along the electron's own direction), the
photon is scattered there with exact Compton energy--momentum conservation,
\begin{equation}
E'_s = \frac{E'}{1+\epsilon\,(1-\cos\theta)},
\end{equation}
and the outgoing electron and photon four-momenta (formed by conservation) are
boosted back to the lab, where the electron momentum is mapped to $(x',y',\delta)$
without approximation.  This matches the analytic Compton edge
$4\gamma^2 E_h/(1+2\epsilon)$ and should be used in the hard inverse-Compton /
$\gamma$-ray regime ($\epsilon\gtrsim 0.1$); at small $\epsilon$ it agrees with the
default to a few percent.  The overlap integral and the sampling of $\theta$ are
identical in both modes.  The default (\verb|EXACT_RECOIL|$=0$) retains the
Eq.~(4) mapping for backward compatibility.

\paragraph{Cross section.}
The scattered-photon polar angle is sampled from the differential cross section
selected by \verb|CROSS_SECTION|:
\begin{itemize}
\item \verb|"thomson"|: $d\sigma/d\Omega \propto (1+\cos^2\theta)$, with total
  cross section $\sigma_T = \frac{8\pi}{3}r_e^2$~\cite{Jackson}.
\item \verb|"klein-nishina"| (default): the full Klein--Nishina
  expression~\cite{KleinNishina1929} with $\epsilon = E'/(m_e c^2)$, and the
  corresponding Klein--Nishina total cross section.
\end{itemize}
In the keV rest-frame-energy regime typical of optical-laser Compton sources the
two choices agree to order $\epsilon$; both are provided for generality.

\paragraph{Scatter probability (luminosity overlap).}
The mean number of scatters for a given particle on a given pass is computed as a
luminosity-overlap integral of that particle's path through the Gaussian laser
pulse,
\begin{equation}
\mu = \sigma_{\rm tot}\,\text{\tt FACTOR}\,\frac{1-\beta\cos\theta_0}{\beta}
      \int n_{\rm ph}\big(\mathbf{r}_e(z),t_e(z)\big)\,dz,
\end{equation}
integrated over the electron beamline coordinate $z$, where $n_{\rm ph}$ is the
laser photon density, $(1-\beta\cos\theta_0)$ is the relative-velocity (flux)
factor --- exact for any angle --- and the $1/\beta$ converts the time integral to
a path-length integral.  The photon number is $N_{\rm ph}=$\verb|PULSE_ENERGY|$/E_h$
unless overridden by \verb|N_PHOTONS|.

The overlap is evaluated in the \emph{laser propagation frame} rather than the
beamline frame, so the geometry is correct for any collision angle, including a
$90^\circ$ side collision.  The propagation direction and two transverse laser
axes are built from \verb|COLLISION_ANGLE| ($\theta_0$, the angle between $+\hat z$
and the laser propagation; $\theta_0=\pi$ is head-on) and \verb|COLLISION_AZIMUTH|
($\phi_0$, measured from $+\hat x$ toward $+\hat y$, so $\phi_0=0$ is a crossing in
the horizontal $x$--$z$ plane and $\phi_0=\pi/2$ a vertical crossing):
\begin{eqnarray}
\hat{\mathbf k} &=& (\sin\theta_0\cos\phi_0,\ \sin\theta_0\sin\phi_0,\ \cos\theta_0)
   \qquad\text{(propagation; \tt LASER\_PULSE\_LENGTH along here)},\\
\hat{\mathbf e}_1 &=& (\cos\theta_0\cos\phi_0,\ \cos\theta_0\sin\phi_0,\ -\sin\theta_0)
   \qquad\text{(in-plane transverse; {\tt SIGMA\_X})},\\
\hat{\mathbf e}_2 &=& (-\sin\phi_0,\ \cos\phi_0,\ 0)
   \hspace{6.2em}\text{(out-of-plane transverse; {\tt SIGMA\_Y})}.
\end{eqnarray}
Thus \verb|SIGMA_X| and \verb|SIGMA_Y| are the laser rms spot sizes along
$\hat{\mathbf e}_1$ (in the crossing plane) and $\hat{\mathbf e}_2$ (perpendicular
to it), and \verb|LASER_PULSE_LENGTH| is the rms length along the propagation
direction $\hat{\mathbf k}$.  At head-on this reduces to $\hat{\mathbf k}=-\hat z$,
$\hat{\mathbf e}_1=\hat x$, $\hat{\mathbf e}_2=\hat y$, i.e.\ the familiar
pulse-length-along-$z$ picture.  Two focusing effects are included simultaneously:
\begin{itemize}
\item {\em Hourglass:} the laser rms spot size grows away from the waist
  \emph{along the propagation direction} as
  $\sigma(s)=\sigma_0\sqrt{1+(s/z_R)^2}$, where $s$ is the distance from the waist
  measured along $\hat{\mathbf k}$, with Rayleigh range $z_R = 4\pi\sigma_0^2/\lambda$
  computed separately for \verb|SIGMA_X| and \verb|SIGMA_Y| from
  \verb|LASER_WAVELENGTH|.  The waist location is set by \verb|FOCUS_POSITION|
  ($S_0$, along $\hat z$) together with the transverse offsets \verb|DX|, \verb|DY|.
\item {\em Crossing geometry:} the electron's separation from the moving pulse
  center is projected onto $(\hat{\mathbf k},\hat{\mathbf e}_1,\hat{\mathbf e}_2)$,
  so both the reduced interaction length of an oblique collision and the exact flux
  factor $(1-\beta\cos\theta_0)/\beta$ are captured.
\end{itemize}
Each particle's true (velocity-corrected) arrival time, referenced to the bunch
centroid and offset by \verb|LASER_DELAY| (a time in seconds), sets its overlap,
so head/tail and off-axis particles sample different, hourglass-broadened overlaps.
Because the exponent is quadratic in $z$, the quadrature is centered on the overlap
peak and spans $\pm 5$ effective sigma about it, evaluated over \verb|N_STEPS|
points (the full hourglass-broadened density is used at each point).  A single
scatter per pass is applied (valid for $\mu\ll 1$); if $\mu>1$ it is clamped to $1$
and a warning is issued.  The overall rate can be scaled with \verb|FACTOR|.

\paragraph{Photon output.}
If \verb|PHOTON_OUTPUT_FILE| is given, the emitted back-scattered photons are
written to an SDDS file with columns \verb|Ep| (photon energy, eV), \verb|x|,
\verb|xp|, \verb|y|, and \verb|yp|.  Photon output is available in the serial
version of {\tt elegant} only.

The output filename may be an incomplete filename.  This means it may contain one instance of the string format
specification ``\%s'' and one occurence of an integer format specification (e.g., ``\%ld'').  {\tt elegant} will replace
the former with the rootname (see \verb|run_setup|) and the latter with the element's occurrence number.
This allows labeling the files from different interaction points with sequential indices.

\paragraph{Notes.}
The element applies at most one Compton event per particle per pass, consistent
with Eq.~(4).  For multi-pass (storage-ring) use, \verb|STARTONPASS| and
\verb|ENDONPASS| gate the passes on which scattering occurs, and
\verb|PASS_INTERVAL| further restricts scattering to every \verb|PASS_INTERVAL|-th
pass counting from \verb|STARTONPASS| (the default of 1 acts on every pass in the
window).

